%----------------------------------------------------------------------------
\chapter{Code, tools and checks}\label{guide:code}

Paper build commands and paper source paths in this chapter are relative to \path{paper/}. Repository layouts and commands beginning with \texttt{python3 scripts/} are marked when they run from the monorepo root. Workflows and agent skills remain at that root.
%----------------------------------------------------------------------------
\section{GitHub}\label{guide:code:github}
Keep the paper and its supporting research in one private project repository (\cref{guide:start:access}). Publish deliberately selected artifacts through Zenodo (\cref{guide:zenodo}); scratch notes are excluded by default. A public research package does not require publishing every file in the working repository.

\paragraph{Repository layout.} Paths here start at the monorepo root:
\begin{itemize}[noitemsep]
  \item \texttt{paper/}: paper, bibliography, figures, guide and paper checks;
  \item \texttt{code/}: implementation and tests, with its Dockerfile and development-container configuration;
  \item \texttt{proof/}: selected formal models and proof-checking profiles;
  \item \texttt{benchmark/}: measurement workloads and reproduction instructions;
  \item \texttt{scratch/}: research anchors, consultation summaries and working notes;
  \item \texttt{presentation/}: presentation sources and rendering tools.
\end{itemize}
Large data, dependencies and generated outputs stay outside tracked source. A release records its selected source revision, accepted input packages, tools, commands and result hashes.

\paragraph{Shared commands.} Root \texttt{project.json} declares the enabled components and their actual commands. From the monorepo root, \texttt{python3 scripts/project.py doctor} inspects prerequisites and \texttt{python3 scripts/ci/run.py} executes the configured ordinary checks with recorded runtime identities. Use a fresh checkout for evidence; repeated runs need new output locations. Unused components report \emph{not configured}. Enabling a component requires real checks; an empty recipe or a successful command that only prints TODO is not evidence.

Put project reproduction commands in a \texttt{justfile} when useful. Implement \texttt{just reproduce} for the paper's results and a reduced smoke workload for routine CI. Name commands after the figures or tables they produce and document the hardware, datasets, seeds and units. The inherited code recipes are unconfigured until you replace them. The top-level benchmark example demonstrates the measurement interface; its timings are not research results.

\section{The builds and what reads them}\label{guide:code:builds}
The template has one set of sources and two \LaTeX{} roots, so it builds two PDFs (\cref{guide:tab:code:builds}). You edit \texttt{draft.tex}, which holds the title, \texttt{\textbackslash lintphase} and the chapter list. \texttt{final.tex} is four lines: it sets the \texttt{final} class option, disables todonotes, defines \texttt{\textbackslash finalmode} and inputs \texttt{draft.tex}. Never edit it. Your text lives in \texttt{content/}, and both PDFs print the same paper from it. The names follow \LaTeX{}'s \texttt{draft} and \texttt{final} options, but only the final build sets its option: the class option \texttt{draft} would replace your figures with empty boxes, so \texttt{draft.tex} leaves it out.
\begin{table}[!ht]
  \centering
  \small
  \caption{The two builds.}
  \label{guide:tab:code:builds}
  \begin{tabularx}{\linewidth}{>{\raggedright\arraybackslash}p{0.17\linewidth} >{\raggedright\arraybackslash}X >{\raggedright\arraybackslash}X}
    \toprule
     & \textbf{Draft} & \textbf{Final} \\
    \midrule
    Root & \texttt{draft.tex} & \texttt{final.tex} \\
    Locally & \texttt{make draft} or \texttt{make all} & \texttt{make final} \\
    On Overleaf & root \texttt{paper-draft.tex} & root \texttt{paper-final.tex} \\
    Output & \texttt{draft.pdf} & \texttt{final.pdf} \\
    Contains & the todo list, the Checks page, the notes, beats and gap markers, the paper, this guide & the paper only \\
    Read by & you and your advisors, while you write & the TDK reviewers, the lint in CI, the self-score and the expert personas \\
    \bottomrule
  \end{tabularx}
\end{table}

\paragraph{\emph{Final} names the format, not the state of the paper.} A final PDF builds from your first day. It hides every todo, beat and gap marker, so it can look finished while half the paper is missing. Read it every week to see what reviewers will see, and run \texttt{make check-submit} to learn whether it is ready. The draft build is also separate from the \texttt{draft} phase of the lifecycle and from the full draft, the anchor we approve (\cref{guide:lifecycle}). You build both PDFs in every phase.

\paragraph{What is made from the builds.} The lint and paper CI produce four more outputs:
\begin{description}[noitemsep]
  \item[Checks page] Part of the draft. On every compile, the lint runs on the sources and prints its findings at the front of \texttt{draft.pdf} (\cref{guide:latex:todo}). It runs no PDF rules and does not read the build log, so a clean Checks page does not promise a green CI.
  \item[Skim view] Made by paper CI from \texttt{final.pdf} (\cref{guide:lifecycle:reviewers}). It keeps headings, labelled items, boxes and captions, and replaces paragraphs with their length.
  \item[Review bundle] An artifact named with the full source SHA, run ID and attempt, kept for 7 days: both PDFs, their logs, the lint's output, the skim view and a manifest (\cref{guide:code:reviewers}). The self-score and the reviewer passes read this bundle, never a PDF on your machine, so what they judge is tied to a commit.
  \item[Burn-down counts] Counts reported by the lint, which can be saved and plotted deliberately (\cref{guide:code:paperci}).
\end{description}

\paragraph{Which check reads which build.} Most rules read only the sources and give the same result for either build. The option \texttt{--doc} picks the build whose log, \texttt{.aux} file and PDF the remaining rules read: \texttt{build}, \texttt{overfull}, \texttt{pages}, \texttt{sources}, \texttt{checkcites}, \texttt{bibtex-log} and every \texttt{pdf-} rule. Read those results on the final build, because the draft's notes and guide change its page count, its margins and its text.
\begin{description}[noitemsep]
  \item[\texttt{make check}] The draft build, at your \texttt{\textbackslash lintphase}. Its PDF rules also measure the notes and this guide, so prefer the next command.
  \item[\texttt{make check-final}] The final build, at your \texttt{\textbackslash lintphase}. Paper CI runs the same rules on the same build.
  \item[\texttt{make check-submit}] The final build, at the \texttt{submit} phase whatever \texttt{\textbackslash lintphase} says. This is the check before you upload.
\end{description}
No reviewer persona reads \texttt{draft.pdf}. The skim reviewer reads the skim view, and the expert personas and the self-score read \texttt{final.pdf}.

\section{The checks and paper CI}\label{guide:code:paperci}
Root \texttt{.github/workflows/ci.yml} checks affected components on ready pull requests, all enabled ordinary components on default-branch pushes, and selected components on manual runs. Draft pull requests receive structure and configuration feedback without claiming full execution. One paper execution builds both PDFs and generates lint and review evidence. Separate \emph{Paper build} and \emph{Paper lint report} workflows read the exact completed run for their badges; they do not rebuild it. A green lint-report badge means the checker produced its report, not that the paper has no findings.

Writing findings are advisory in ordinary CI. Broken compilation, checker failures, invalid configuration and missing required evidence fail. Before submission, use the strict submission check without changing the student's phase. There are three rule families; \texttt{texlua scripts/lint.lua list} describes each rule, with its severity and phase.
\begin{description}[noitemsep]
  \item[source] Checks the sources and build log. Gates cover phase-appropriate essentials (\texttt{essentials}), the template structure (\texttt{frame}) and a clean build (\texttt{build}, \texttt{overfull}); the chapter-order check (\texttt{spine}) only advises. They also check floats, captions, axes, tables, labels, references, acronyms, listings, the shape of the reference list, typography, consistent terms, that no TDK paper of the same year is cited or linked (\cref{guide:tdk:companion}), headline numbers against the evaluation, and completion of the AI declaration and title page. The burn-down rules for placeholders and beats only report, and the \texttt{burndown} rule writes their counts as one line of JSON for CI to collect over time. A rule that cannot run, because a tool is missing or the profile excludes it, says \emph{skipped}, never \emph{pass}.
  \item[pdf] Checks the final PDF against the form requirements (\cref{guide:tdk:form}), even if you replaced the template. Checks cover the paper size, metadata, embedded fonts, a 15\,MB limit and the body font: \texttt{pdf-font} fails on a font that is not a serif or not 12\,pt, and only warns on a serif other than Times. The 15\,MB limit is our safeguard, taken from the faculty's thesis guideline (\texttt{include/guideline.tex}); it is not a TDK rule. They measure line spacing from line pitch and margins from the text on every page. They report text below 10\,pt, math scripts aside, count main-text pages from the PDF outline, and check required parts, including both abstracts and the AI declaration. For a thesis, \texttt{LINT\_PROFILE=thesis} applies the faculty's thesis thresholds.
  \item[tools] Reports from TeX Live and container tools: chktex and lacheck for typography, texcount for chapter word counts and abstract length, checkcites for uncited reading-list entries, the BibTeX log, aspell on the detexed chapters for unknown words and mixed British and American spellings, and a link check. These report and never block the build, with one exception: \texttt{repo-hygiene} fails at every phase when the repository tracks build output such as a PDF, a log or an \texttt{.aux} file.
\end{description}
The \texttt{quotations} rule requires a citation in each \texttt{sourcequote} attribution from \texttt{draft} onward. It also warns about missing locators and suggests the styled environment for plain \texttt{quote} or \texttt{quotation} blocks. The \texttt{bib} rule checks that citation keys exist. The lint does not require a citation after every \texttt{\textbackslash enquote} (which also marks terms and invented examples), or in every paragraph. Check inline quotation attribution, whether a source supports a claim, and whether attribution remains clear across paragraphs yourself (\cref{guide:writing:citation-placement}). It cannot verify whether \texttt{[sic]}, bracketed replacements or omissions faithfully represent the original.
Each rule has a \emph{severity} and a \emph{phase}. A \emph{check} fails the run, \emph{advice} warns, and a \emph{count} reports numbers. A check whose phase is \texttt{always} fails at every phase: \texttt{essentials}, \texttt{frame}, \texttt{guide:chapters}, \texttt{appendix-core}, \texttt{build}, \texttt{overfull}, \texttt{figures}, \texttt{labels}, \texttt{crefs}, \texttt{bib}, \texttt{companion}, \texttt{code}, \texttt{preamble} and \texttt{repo-hygiene}. The other checks become required at \texttt{draft}, \texttt{polish} or \texttt{submit}; before then, they only report. The phases are those of the lifecycle (\cref{guide:lifecycle}): \texttt{topic}, \texttt{setup}, \texttt{research}, \texttt{plan}, \texttt{draft}, \texttt{polish} or \texttt{submit}. The first four share the checks' first level, and the lint rejects any other value. In CI, a raw lint exit of 1 reports an advisory paper finding; raw exits above 1 and other technical errors fail the workflow. The template repository also runs regression and consistency checks. Student copies do not run template-skeleton checks. Locally, \texttt{make check} runs every rule against the draft build, \texttt{make check-final} against the final PDF, and \texttt{make check-submit} against the final PDF at the \texttt{submit} phase whatever \texttt{\textbackslash lintphase} says, the check before you upload; run one rule with, for example, \texttt{make lint-essentials}. The same lint generates the draft PDF's \emph{Checks} page on every compile (\cref{guide:latex:todo}), keeping the PDF and CI consistent.

The advisor creates the GitHub project before importing it into Overleaf. View its runs under GitHub's \emph{Actions} tab or on a pull request. CI records the source SHA, exact configuration, resolved tool versions and file hashes. Ordinary evidence lasts 14 days; bulky previews and review bundles last 7 days. Save evidence needed for a publication in its durable release package before it expires.

\paragraph{Filled evaluation results.} Run \texttt{make lint-evaluation-results} to check literal measurement columns (such as Time, Latency, Throughput, Memory or Accuracy) and explicit \texttt{\textbackslash num}, \texttt{\textbackslash SI} and \texttt{\textbackslash qty} values in evaluation chapters. Blank cells, placeholders and unexplained dashes or \texttt{N/A} are findings; zero is a valid value, and \texttt{N/A: <reason>} explains an unavailable measurement. The rule reports during drafting and becomes required at \texttt{polish}; ordinary CI still treats writing findings as advisory. It reads simple \texttt{tabular}, \texttt{tabular*}, \texttt{tabularx} and \texttt{longtable} tables, excluding notes, comments and code. Complex tables, macro-generated results, plots and unmarked prose need human review. This checks whether values are filled, not whether they are correct or supported by experiments.

\paragraph{Official references.} Ordinary checks validate accepted local reference packages. Weekly maintainer checks report changed or inaccessible official sources and overdue applicability review. A network outage does not invalidate verified cached bytes. Missing or corrupt required bytes fail the dependent review. The Overleaf Checks page cannot claim artifact-dependent rubric verification when those reference bytes are unavailable.

\paragraph{Lifecycle reports.} Run \texttt{python3 scratch/scripts/lifecycle.py} from the monorepo root when you want to save a lifecycle inventory. It reads the paper phase, research anchors and consultation summaries, and component-specific Git history. Reminders about missing or stale research work are advisory. They do not change the phase or automatically commit generated notes. These checks inventory the work; we still assess its quality.

The repository move in the following retained advisor note has been completed by the monorepo integration.
\advisortodo{Move \texttt{tdk-notes-template} and \texttt{tdk-code-template} to an FTSRG organisation when there is one for the lab; \texttt{ftsrg-bta} is the course's.}

\section{Running the reviewer passes}\label{guide:code:reviewers}
\paragraph{Self-score and reviewer passes.} The monorepo carries two agent skills in root \path{.claude/skills/}, so each commit has the skills that review it. On main-branch pushes, non-draft pull requests and manual runs, paper CI uploads the review bundle of that commit (\cref{guide:code:builds}). \path{.claude/skills/self-score/} scores the artifact's \texttt{final.pdf} A--E using the review sheet's wording and page references, and identifies the weakest criterion and cheapest fix. \path{.claude/skills/reviewer-passes/} runs the skim reviewer and expert panel on the same artifact. \texttt{scripts/review/review.py} fetches the artifact of a named commit and checks its hashes. Self-score writes nothing; the reviewer passes write nothing unless you record a run. \Cref{guide:lifecycle:reviewers} says what the passes are and when to run them; this section says how.

\paragraph{Running them.} Push your commit and run paper CI, through a non-draft pull request or a manual workflow dispatch (\cref{guide:code:paperci}). Then, in your paper repository, ask your agent to run the reviewer passes, or name one persona: \enquote{review as the cryptographer}. The skill in \path{.claude/skills/reviewer-passes/} fetches the CI artifact of that commit with \texttt{scripts/review/review.py}, which refuses an artifact whose hashes do not match. It gives each persona a directory holding only that persona's inputs: the final PDF, the skim view, or both, and your spine if you name the file. Without a spine, a persona judges overclaiming against the abstract and the introduction. Each persona runs apart from the others and never sees another persona's findings or an earlier run. The chat response gives each persona's worst finding, the criterion and page, and the cheapest fix. It then lists the three cheapest changes overall. Ask a persona for more to see its remaining findings.

\paragraph{Recording a run.} By default the skill writes no files and never edits the paper. You choose the fixes, and the consultation summary (\cref{guide:consult}) records your decisions. We record a run at three milestones (the spine premortem, the full draft and before upload), and you may record any run you choose. A recorded run becomes one issue on your paper repository: the panel, the skim scores, one checkbox per finding, and the hashes of the artifact and of each persona's inputs. Propose a disposition for each finding in a comment on the issue: \texttt{F3 fixed in <sha>}, \texttt{F4 rejected: <reason>} or \texttt{F5 deferred}. We confirm the dispositions with the label \texttt{review-confirmed}. The issue is not an input to the next run, and the skill keeps no memory between runs. Rerun it on the revised paper to check whether the fixes answer the objections.

Ignore findings that do not identify a place in the paper. Verify every suggested source before citing it; a persona may invent one. Tell us if two personas raise the same concern or one argues outside its remit, so we can fix their boundaries in the persona files.

\section{The README}\label{guide:code:readme}
Every TDK needs reproduction instructions for its results. Some also produce a reusable \emph{tool}. Expand \texttt{code/README.md} and \texttt{benchmark/README.md} with the following sections as the project develops. Do not claim a release, result or DOI before it exists.

\paragraph{The results README.} In this order:
\begin{enumerate}[noitemsep]
  \item What this is, in two sentences, with the paper's title and the concept DOI badge (\cref{guide:zenodo}).
  \item The reproduction command: \texttt{docker run \dots{} reproduce} for every figure and table (\cref{guide:code:docker}). Follow it with \texttt{just reproduce} for local clones and the exact dependency versions resolved for that run. Keep dependencies to a minimum.
  \item A results map with one row per paper figure and table: its number, reproduction command, output path and wall-clock time on specified hardware.
  \item Reproduction limits: secrets, special hardware, long runs or nondeterminism, including the seed policy and expected variance. Keep this section even if everything reproduces with fixed seeds.
  \item The results tag: which Git tag and Zenodo version the paper's numbers came from (\cref{guide:zenodo:versions}).
  \item How to cite, and the license.
\end{enumerate}

\paragraph{The tool README} adds installation, use on the reader's own input with a worked example, contribution instructions and maintainers after the results sections. Follow the conventions of the open-source foundations we work with, such as \href{https://github.com/LF-Decentralized-Trust/lab-proposals/blob/main/README.md}{LF Decentralized Trust Labs}: an Apache 2.0 \texttt{LICENSE}, a \href{https://developercertificate.org/}{Developer Certificate of Origin} sign-off on every commit (\texttt{git commit -s}), \texttt{MAINTAINERS.md} and a \texttt{SECURITY.md} with vulnerability-reporting instructions. A compliant tool can be proposed as a lab after the TDK.

\section{Docker}\label{guide:code:docker}
Use Docker to make the research environment available to a reviewer. The starter's \texttt{code/Dockerfile} builds a development scaffold with current packages and Just; its default command lists the still-unconfigured recipes. The code development-container configuration uses that same Dockerfile. It does not yet build the project's implementation or reproduce a result.

Once implemented, document a container command that regenerates the selected results into a mounted output directory. Normal CI resolves latest tools once per run and records their actual versions and image digests. Preserve the exact image identity and input/output hashes with published evidence; mutable GHCR tags do not freeze image bytes. Choose explicit publication commands for a working research image. The retained nested legacy workflows are inactive; root CI is the shared execution interface.

\section{Demos}\label{guide:code:demo}
Record demonstrations with a script in the code repository, never by hand. A script lets you update the recording after code changes and rerun it before the presentation. Configure a smoke recording in ordinary CI when affordable; run longer demonstrations explicitly. Save the selected recording with its source and runtime evidence for reviewers.
\begin{itemize}[noitemsep]
  \item For a web UI, we recommend \href{https://playwright.dev/}{Playwright}: it drives the real workflow and saves a video and trace. Draw captions into the frames at each step.
  \item For a backend, CLI or prover such as Tamarin, Lean or a model checker, we recommend recording a scripted terminal session with \href{https://asciinema.org/}{asciinema}. Render it to SVG or GIF, or embed the cast beside the verdicts in HTML. We use this for our backends and proof assistants.
\end{itemize}
Use captions to explain the invariant, rejected attack or result in the running example's terms. Do more than narrate clicks. The starter has demonstration guidance under \texttt{code/demo/}. Implement the recorder and \texttt{just demo} before enabling its checks.

\section{Behaviour-driven tests}
For behaviour-driven development (BDD), you can write the expected behaviour as scenarios with \href{https://cucumber.io/}{Cucumber}; the scenarios then double as evaluation cases that a reviewer can read.

\section{AI tools}\label{guide:code:ai}
\paragraph{Tools, as of October 2026.} These recommendations can become outdated within months.
\begin{itemize}
  \item For research, we recommend the strongest models. We prefer Claude Fable and ChatGPT Astra; ask whether we can help with a subscription.
  \item Prefer an agent working in your repository: Claude Code or Codex, ideally through T3 Code, GitKraken Kepler or herdr. This gives the agent access to your files and its repository tools. You \emph{need} version control when AI edits files.
  \item Plugins and skills: the Matt Pocock and poetato plugins for productivity; caveman, ponytail, squeez and TRON if you want to spare tokens. Three we use daily: \href{https://www.aihero.dev/skills-grill-me}{grill-me} (\href{https://claudemarketplaces.com/skills/mattpocock/skills/grill-me}{marketplace page}) structures a Socratic dialogue with the AI (\cref{guide:code:grill}); \href{https://github.com/cursor/plugins/blob/main/pstack/skills/unslop/SKILL.md}{unslop} takes the machine out of generated prose; \href{https://github.com/DietrichGebert/ponytail/blob/main/skills/ponytail/SKILL.md}{ponytail} cuts the tokens an agent spends.
  \item We recommend concise instructions. Overloaded \texttt{AGENTS.md} files and memories can impair model performance. Prefer skills that clearly describe the range of tasks they apply to.
\end{itemize}

\subsection{Socratic dialogue with AI}\label{guide:code:grill}
Use a Socratic dialogue with the AI to work through your understanding: you explain a claim or decision, the AI asks what it means and what supports it, and you test and refine the answer together. Ask it to expose assumptions, request examples and propose counterexamples. You should leave able to reconstruct the reasoning yourself. Challenge the AI's premise when it does not fit your work, and check its answers as carefully as it checks yours.

The \texttt{grill-me} skill is a way to structure this dialogue. Use it on a bounded question: the problem in your topic manifest, a contribution in your spine, an evaluation design or a concept you cannot yet explain. Give the agent the relevant artifact and say what you want to understand or settle. A whole thesis is too broad for one useful round.

\paragraph{An example.} Suppose you tell the AI that a signed diploma status proves the diploma is still valid. It asks what the signature establishes. You answer that the issuer signed the status. It then asks whether a verifier could receive an authentic status after the diploma was revoked. Your answer exposes the need for a freshness condition. The dialogue continues: how old may the evidence be, whose clock measures its age, and what happens when no recent evidence is available? Each answer changes what the next question can assume. You leave with a narrower claim, explicit timing assumptions and a test with delayed evidence, or a specific question that needs a proof or a source.

\paragraph{Make the answer yours.} Explain in your own words, with a small example if needed. If you cannot explain a term or an AI-generated argument, ask for a hint or explanation, then apply it to a different case. Agreement alone does not show understanding; predicting what happens in the changed case does.

Treat the decisions as a tree. First settle the choices on which other choices depend. For example, agree on the intended use and acceptable evidence age before choosing refresh intervals or benchmark workloads. Ask questions whose prerequisites are settled in the current round, answer them, then move to the questions those answers unlock. The agent should inspect available files for facts; you supply your intended scope and make the decisions. If evidence is missing, record the investigation needed rather than accepting a guess.

\paragraph{A prompt to adapt.}
\begin{quote}
Use grill-me to help me test the evaluation plan in [file]. Read the plan and the claims it is meant to check. Ask questions in dependency order, waiting for my answers between rounds. Challenge vague terms, unsupported assumptions and cases where the proposed evidence would not establish the claim. Look up facts available in the repository yourself and identify anything you cannot verify. For each decision, explain your recommendation and its tradeoffs; ask me to justify my choice and test it on a counterexample. Do not edit the plan during the interview. Finish by summarising the decisions, their reasons, unresolved questions and the evidence I need next. Ask me to confirm that the summary matches my understanding.
\end{quote}

\paragraph{Use recommendations critically.} A fluent recommendation can make agreement too easy. Explain why it fits your case and what would make you reject it. Ask the agent to challenge that explanation. If the discussion becomes repetitive, ask which unresolved decision blocks progress and what evidence would resolve it. The useful output is a decision you can defend, or a specific investigation to carry out.

End the dialogue when you can explain the decision and its reasons, or when further progress needs evidence you do not yet have. Record what would settle each remaining question and the next investigation to carry out. Save the decisions and open questions in your notes and retain the conversation as described in \cref{guide:writing:ai}. Bring the artifact and the questions that remain to consultation. An AI interview does not approve an anchor: changes to the topic manifest, spine or approved draft still need our agreement (\cref{guide:lifecycle}). Verify suggested sources before citing them and account for AI assistance in your declaration.
